This section records the manuscript's conditional reduction step in explicit dependency form.
The purpose is to isolate the final logical deduction from the earlier bridges that still require proof.

\subsection*{Conditional reduction statement}
\begin{proposition}[Framework-internal conditional RH reduction]
Assume the manuscript definitions of the state space $H$, admissibility class, and balance criterion are fixed and well-posed under a hypothesis package $\mathcal{A}$.
Assume further that the following three inputs hold under the same $\mathcal{A}$:
\begin{enumerate}
  \item \textbf{Target A (zero--balance correspondence).} Every nontrivial zero in the critical strip corresponds to a balanced point in the admissible class (with converse direction as specified in the correspondence theorem).
  \item \textbf{Target B (critical-line localization of balanced points).} Every admissible balanced point satisfies $\RePart(s)=\tfrac12$.
  \item \textbf{Global off-line exclusion package.} No admissible off-line balanced configuration exists in the full critical strip, including boundary/degenerate regimes required by the framework's parameterization.
\end{enumerate}
Then every nontrivial zero $\rho$ of $\zeta$ satisfies $\RePart(\rho)=\tfrac12$.
\end{proposition}

\subsection*{Separation of burdens: local uniqueness vs global exclusion}
For audit clarity, the manuscript treats two logically distinct obligations.
\begin{enumerate}
  \item \textbf{Local uniqueness-of-balance burden (Target B core).}
  Prove that in a specified regime $\mathcal{R}_{\mathrm{loc}}\subseteq\mathcal{A}$, balance implies $\RePart(s)=\tfrac12$.
  This is a localization statement.
  \item \textbf{Global off-line exclusion burden (Stage 6 core).}
  Prove that no admissible balanced point with $\RePart(s)\neq\tfrac12$ exists anywhere in the full strip region of interest, including non-perturbative, boundary-near, and large-parameter regimes not automatically covered by $\mathcal{R}_{\mathrm{loc}}$.
\end{enumerate}
The second item is not a restatement of the first.
It is an extension/closure problem requiring additional uniform control.

\subsection*{Why symmetry alone is insufficient}
Classical symmetries give transport rules ($s\mapsto1-s$, $s\mapsto\bar s$).
They do not by themselves imply that the only balance-supporting locus is $\RePart(s)=\tfrac12$.
An argument can preserve symmetry and still allow paired off-line candidates.
Hence the framework requires a separate exclusion mechanism beyond symmetry compatibility.

\subsection*{Reduction proof skeleton (conditional)}
The deduction itself is short once the above inputs are available.
Let $\rho$ be a nontrivial zero.
By Target A, $\rho$ maps to a balanced admissible configuration.
By Target B, that balanced configuration must lie on $\RePart(s)=\tfrac12$.
Hence $\RePart(\rho)=\tfrac12$.
The global off-line exclusion package ensures that no admissible off-line cases survive outside the local regime used to prove Target B, so the conclusion applies across the entire critical strip addressed by $\mathcal{A}$.

\subsection*{Where the argument is formally straightforward}
\begin{itemize}
  \item The final implication from ``zero is balanced'' plus ``balanced implies critical line'' to RH-in-strip is a direct composition of statements.
  \item No deep analytic estimate enters at this final step once hypotheses are synchronized and side conditions are already checked.
\end{itemize}

\subsection*{Where the unresolved burden lies}
The mathematical pressure is upstream, not in the final deduction.
The unresolved burden is to prove, under explicit assumptions, the six core open bridges that feed the reduction:
\begin{enumerate}
  \item \textbf{Exact spectral object burden:} finalize the $H$-layer with explicit domain/action/symmetry data and proof-level well-posedness.
  \item \textbf{Zero--balance correspondence burden (Target A):} establish a rigorous correspondence map with precise domain/codomain control, including both directions and exclusion of spurious balanced states.
  \item \textbf{Uniqueness-of-axis burden (Target B):} prove a theorem-level localization statement forcing every admissible balanced point onto $\RePart(s)=\tfrac12$.
  \item \textbf{Off-critical-line exclusion burden:} upgrade local or regime-restricted arguments to a full-strip exclusion theorem compatible with the assumptions used in Targets A and B.
  \item \textbf{Operator-theoretic burden:} prove each operator property actually used in Targets A/B and exclusion (no hidden assumptions).
  \item \textbf{Analytic-control burden:} justify every growth/convergence/interchange step needed to pass from local constructions to strip-wide statements.
\end{enumerate}
Until these bridges are proved, the reduction remains conditional.

\subsection*{Insufficient argument patterns (explicitly excluded)}
The manuscript does not treat the following as adequate for Target~B or global exclusion:
\begin{enumerate}
  \item \textbf{Pure symmetry arguments} that establish only mirrored pairing of candidate points.
  \item \textbf{Local-only coercivity} without a theorem transferring estimates to the full admissible strip.
  \item \textbf{Regime-fragmented proofs} whose hypotheses differ across subregions without compatibility/patching lemmas.
  \item \textbf{Numerical-only exclusion} not converted into theorem-level analytic statements.
\end{enumerate}

\subsection*{Assumption-compatibility requirements}
For the reduction to be valid as stated, the assumptions used in Target A, Target B, and global off-line exclusion must be mutually compatible.
In particular, the manuscript must avoid: (i) proving Target A on one admissible class and Target B on a narrower incompatible class without a transfer theorem, and (ii) invoking a global exclusion principle whose hypotheses are not inherited by the correspondence setup.

\begin{remark}[Audit-facing interpretation]
This section does \emph{not} claim that RH is proved.
It states that the framework yields RH \emph{if} Targets A and B and the global off-line exclusion package are established under a common explicit hypothesis set.
The hard work is proving those prior bridges.
\end{remark}

\begin{remark}[Common failure modes in RH-type reductions]
Two recurring failure modes are especially relevant here.
First, a correspondence theorem may hold only with hidden regularity assumptions, leaving uncovered zero candidates.
Second, a localization argument may be valid only in a perturbative or local regime and fail to imply global off-line exclusion.
The present framework therefore keeps Target A, Target B, and global exclusion as separate proof obligations.
\end{remark}

\openbridge{Current manuscript status: the conditional deduction is structurally clear, but the six core bridges (exact spectral object, zero--balance correspondence, uniqueness of the balance axis, off-critical-line exclusion, operator-theoretic closure, analytic-control closure) remain open proof obligations requiring full theorem-level proofs under synchronized assumptions.}
